\scanpage{639}
\mspage{27}

(\unsure{2.7.1}) $\deleted{\unsure{\mathrm{Mono}}}C\subset\operatorname{Cof}_W$, i.e. \deleted{pour \unsure{tout|tt} diag}\inserted{la condition}

\unsure{a'}) Pour tt diagramme cocartésien

\diagram{\begin{tikzcd}[ampersand replacement=\&,row sep=small]
A_0 \arrow[r,"i",no head] \arrow[d,"f"] \& A_1 \arrow[d,"f'"]\\
A_2 \arrow[r,"i'",no head] \& A_3
\end{tikzcd}}

avec $i$ mono, si $f\in W$, alors $f'\in W$.
(Stabilité de $W$ par \unsure{cochangemt} de base
par monomorphismes.)

Sous ces conditions, ce triple de Quillen est
spécial, et a fortiori clos.

Ce corollaire résulte de 2.6. $3^\circ$), puisque
$C\subset\operatorname{Cof}_W$ implique $TC\overset{\mathrm{def}}{=}C\cap W=C\cap W^{\unsure{\mathrm{univ}|\mathrm{uni}}}$,
\deleted{donc que} (car $\operatorname{Cof}_W\cap W=W^{\mathrm{univ}}$), donc que
$TC$ est stable par \unsure{coch.} de base, d'où la
condition (2.6.5 a).

2.8. On dit qu'une \uline{structure de Quillen}
$(W,C,F)$ sur $M$ est \uline{très} \uline{spéciale}, si elle
est obtenue comme au th. 2.1, i.e.
si on a les conditions

\[
\left\{\begin{aligned}
\text{a)}\quad & M\ \text{est très spéciale, i.e. satisfait aux}\\
 &\text{conditions (2.6.1)}\\
\text{b)}\quad & W\ \text{est stable par }\inserted{\text{petites}}\ \varinjlim\ \text{filtrantes et}\\
 &\text{est accessible (conditions (2.6.3))}\\
\text{c)}\quad & C=\mathrm{Mono}.
\end{aligned}\right.\tag*{(2.8.1)}
\]

Donc pour qu'un localiseur $W$ \unsure{de|dans} $M$ corresponde
à une structure de Quillen \deleted{spé} très spéciale
% archive sheet number: 638
